\documentclass[12pt]{article}
\usepackage[T1]{fontenc}
\usepackage[utf8]{inputenc}
\usepackage{lmodern}
\usepackage{textcomp}
\usepackage{enumitem}
\usepackage{geometry}
\usepackage{graphicx}
\usepackage{amsmath, amssymb}
\geometry{margin=1in}
\begin{document}
\begin{center}
Sociology - Graduate Level Assessment 9 \
Total Marks: 40 \
Answer all questions.
\end{center}

\section*{Multiple Choice (10 marks)}
\begin{enumerate}[label=\arabic*.]
\item Sociology differs from common sense in that:
\begin{enumerate}[label=(\alph*)]
    \item it focuses on the researchers' own experiences
    \item it makes little distinction between the way the world is and the way it ought to be
    \item its knowledge is accumulated from many different research contexts
    \item it is subjective and biased
\end{enumerate}
\item Which of the following is not a 'research purpose'?
\begin{enumerate}[label=(\alph*)]
    \item triangulation
    \item explanation
    \item description
    \item exploration
\end{enumerate}
\item An ecclesia is:
\begin{enumerate}[label=(\alph*)]
    \item a religious organization that claims total spiritual authority over its members
    \item a church organized around voluntary rather than compulsory membership
    \item a sect or cult with a very small following
    \item a hierarchy of priests or other spiritual leaders
\end{enumerate}
\item In Durkheim's work, the term 'collective representations' refers to:
\begin{enumerate}[label=(\alph*)]
    \item effervescent ceremonies that create a feeling of belonging
    \item images of gods or totems that are widely recognized
    \item shared ideas and moral values, often symbolized by an object or figurehead
    \item ideological tools used to obscure class divisions
\end{enumerate}
\item The terms 'crisis of the 1970 s' is used to refer to:
\begin{enumerate}[label=(\alph*)]
    \item declining profits and rising unemployment
    \item the eradication of the welfare state
    \item rising divorce rates and the decline of the traditional family
    \item an unfortunate twist in fashion sensibility
\end{enumerate}
\end{enumerate}

\section*{True / False (10 marks)}
\textit{Answer True or False}
\begin{enumerate}[label=\arabic*.]
\setcounter{enumi}{5}
\item True or False: The experimental method relies on using personal beliefs and values to determine the focus of research studies.
\item True or False: Religious organizations such as the Church of Norway and Islam are characterized by their monotheistic beliefs.
\item True or False: Patterson's study of Brixton revealed that black and white residents competed for economic resources and housing.
\item True or False: In stage 3 of the health transition, acute, infectious diseases like typhus and cholera remain the primary causes of illness and death.
\item True or False: A moral panic arises when audiences challenge the ethnic stereotypes represented in media, prompting societal concern.
\end{enumerate}

\section*{Long Form Response (20 marks)}
\textit{Answer each question with a clear and complete explanation.}
\begin{enumerate}[label=\arabic*.]
\setcounter{enumi}{10}
\item Discuss the implications of a population pyramid that is wider at the bottom than at the top, considering its relation to birth and death rates over several generations.
\item Discuss the socio-economic policies of Thatcher's government, focusing on its priorities and the implications of its approach to welfare for single parents, students, and the unemployed.
\end{enumerate}

\vfill
\noindent\textit{End of Paper}
\end{document}